\documentclass[10pt]{ctexart}
\usepackage[a4paper,margin=0.9in]{geometry}
\usepackage{amsmath,amssymb,booktabs,array,tabularx}
\usepackage{enumitem}
\usepackage{listings}
\usepackage{xcolor}
\usepackage{microtype}
\usepackage{hyperref}
\usepackage{xurl}

\setlist[itemize]{leftmargin=1.5em}
\setlist[enumerate]{leftmargin=1.8em}
\setlength{\parskip}{0.35em}
\setlength{\parindent}{0em}
\renewcommand{\arraystretch}{1.15}

\lstset{
  basicstyle=\ttfamily\small,
  breaklines=true,
  columns=fullflexible,
  keepspaces=true,
  frame=single,
  framerule=0.2pt,
  xleftmargin=0.5em,
  xrightmargin=0.5em
}

\hypersetup{
  colorlinks=true,
  linkcolor=blue,
  citecolor=blue,
  urlcolor=blue
}

\title{023：SSCA Summary 质量审计与 Prompt 改造建议\\从 Full Trajectory Dump 到 Compact Observation Trace}
\author{}
\date{2026-06-26}

\begin{document}
\maketitle
\vspace{-1.2em}

\section*{摘要}
本文基于最新 ALFWorld SSCA retry smoke trace，对第一版 summary action 的质量进行人工审计，并给出下一版 prompt 与轨迹压缩方案。结论是：当前实现的 \texttt{full trajectory} summary prompt 可以跑通
\begin{lstlisting}
traj1 -> summary -> reset -> traj2
\end{lstlisting}
的数据链路，但它把完整 prompt、history wrapper、admissible actions 和截断碎片都交给 summary 模型，信息噪声过大。结果是 summary 经常写成“结构化反思”，但未必忠实、可执行，也未必能改善 retry policy。

建议下一版把 full trajectory dump 改成 \textbf{compact observation trace}：仿照 ALFWorld agent 每一步收到的 observation 方式，只保留 task、initial scene、每步环境 observation、projected action、valid flag、environment feedback 和 outcome。summary prompt 则从开放式 policy note 改成 evidence-grounded retry memory，并配套 summary quality gate 与 ablation。

\section{Trace 审计结论}
审计对象是：
\begin{lstlisting}
EXPS/analysis/022_ssca_retry_smoke/
  ssca_retry_grpo_qwen25_15b_alfworld_smoke_seed2026_ssca_review_20260626_135719/
\end{lstlisting}

该 trace 含 8 个 task index，原始 96 条 rows，按 \texttt{(index, phase, step, traj\_uid, response)} 去重后 72 条。由于 smoke 很短，几乎所有 reward 都是 0，所以本次审计重点不是判断算法收益，而是判断 summary 是否满足三个条件：
\begin{enumerate}
  \item \textbf{忠实性}：summary 是否只引用 trajectory 中真实出现的信息；
  \item \textbf{可执行性}：summary 是否给出 ALFWorld 风格的具体 retry 策略；
  \item \textbf{约束力}：traj2 是否减少 malformed / invalid / repeated actions。
\end{enumerate}

\subsection{逐项观察}
\begin{center}
\small
\begin{tabularx}{\linewidth}{c l l X}
\toprule
Index & Task & 质量判断 & 主要问题 \\
\midrule
0 & find two newspaper and put them in sofa & 差 &
结构完整但空泛；误把 ``Welcome to TextWorld'' 当成问题；没有明确已查位置和下一步搜索优先级；traj2 仍有 3/4 invalid。 \\
1 & find two newspaper and put them in sofa & 差 &
编造 \texttt{coin}；失败归因不准确；retry 输出 \texttt{[Action]}、markdown、缺标准 tag。 \\
2 & put peppershaker on drawer & 中等偏差 &
抓住 \texttt{countertop 1} 可能有 peppershaker，但错误描述第一轮已查多个位置；traj2 退化为连续 \texttt{go to countertop 1}。 \\
3 & put peppershaker on drawer & 差 &
幻觉 ``起始观察 missed countertop''；对 countertop 是否有目标物的判断不稳定；retry 仍重复动作。 \\
4 & examine newspaper with desklamp & 差 &
把 desklamp 当 receptacle；建议继续 diningtable，但 traj1 已显示 diningtable 上没有 newspaper；traj2 全部重复去 diningtable。 \\
5 & examine newspaper with desklamp & 很差 &
编造 \texttt{Soapbox}、\texttt{Desk}；失败原因和真实 traj1 不一致；retry 3/4 invalid。 \\
6 & put candle in toilet & 差 &
把 valid 的 \texttt{go to toilet 1} 说成无效；没有指出必须先找 candle；retry 出现 malformed \texttt{[think]}。 \\
7 & put candle in toilet & 中等偏差 &
能指出 drawer1 没有 candle 和重复 open，但没有给出具体搜索顺序；retry 仍重复 drawer/toilet。 \\
\bottomrule
\end{tabularx}
\end{center}

\subsection{总体模式}
这些 failure pattern 指向同一个问题：当前 summary prompt 过于开放，而且输入不是“环境轨迹”，而是“prompt 日志”。模型容易被重复的 task instruction、history wrapper、admissible action list 和截断片段干扰，于是写出泛化反思或编造状态。

这类 summary 若直接作为 actor response 进入 \texttt{uid=:retry} 组训练，会带来两个风险：
\begin{itemize}
  \item 模型学会 verbose reflection，而不是有效 retry memory；
  \item 由于短 smoke 大多 reward 为 0，坏 summary 和好 summary 没有足够 outcome 信号区分，训练会把噪声 summary 也当作可接受行为。
\end{itemize}

\section{问题根因：Full Trajectory Dump 噪声过大}
当前 summary prompt 中的 \texttt{[First trajectory]} 基本记录每一步的 \texttt{obs["text"]}。这对可追溯性有帮助，但不适合直接作为模型输入。原因包括：
\begin{enumerate}
  \item 每一步 observation 里重复包含 ALFWorld role instruction、task、history 和 admissible action 列表；
  \item history 中又嵌套过去的 observation/action，导致上下文递归膨胀；
  \item prompt 被字符级截断后，会出现 ``sponding actions'' 这类碎片，诱导 summary 误读；
  \item 大量 admissible actions 会稀释真正重要的环境变化；
  \item raw model response 中的 malformed 内容未被结构化标注，summary 很难分辨是环境失败还是格式失败。
\end{enumerate}

因此，下一版不应把完整 prompt trajectory 直接交给 summary。更合理的方式是仿照 ALFWorld 自己给 agent 的 state 表达方式，把第一轮轨迹压缩成 environment-centric trace。

\section{建议一：使用 Compact Observation Trace}
建议新增配置：
\begin{lstlisting}
+ssca.summary_context_mode=compact_obs_trace
\end{lstlisting}

核心格式如下：
\begin{lstlisting}
[Task]
find two newspaper and put them in sofa.

[Initial scene]
You are in the middle of a room. Looking quickly around you, you see ...

[Compact first attempt]
Step 1
Obs: <raw environment observation / anchor_obs only>
Projected action: inventory
Valid: true
Env feedback: You are not carrying anything.

Step 2
Obs: You are not carrying anything.
Projected action: look
Valid: false
Malformed reason: contains Chinese / missing action tag / not one exact admissible action
Admissible action excerpt: go to coffeetable 1; go to drawer 1; ...
Env feedback: You are in the middle of a room. Looking quickly around you, you see nothing.

[Outcome]
Reward: 0.0
Won: false
Invalid action count: 3
Final observation: You arrive at sofa 1. On the sofa 1, you see ...
\end{lstlisting}

实现原则：
\begin{itemize}
  \item task 和 initial scene 只出现一次；
  \item 每步只保留 \texttt{Obs -> Projected action -> Valid -> Env feedback}；
  \item 默认不保存完整 admissible actions，只在 invalid 时保存当步 admissible action excerpt；
  \item raw model response 默认省略或截断，只有 malformed 时保留，并附上 malformed reason；
  \item 长 episode 只保留最近 $k$ 步、invalid 步、首次看到目标物的步骤、目标相关动作步骤和 final observation。
\end{itemize}

这样 summary 看到的是状态变化，而不是 prompt 渲染细节。它更接近 ALFWorld agent 实际面对的 state，也更利于后续迁移到 WebShop / AppWorld 等环境。

\section{建议二：改成 Evidence-Grounded Retry Memory Prompt}
当前 \texttt{SUMMARY\_INSTRUCTION} 太像开放式反思。建议把目标改成更硬的 evidence-grounded schema：
\begin{lstlisting}
You are writing an evidence-grounded retry memory for the same ALFWorld task.
Use only facts supported by the compact first attempt.
Do not invent objects, receptacles, locations, or actions.

Return exactly these fields:

Task:
Observed facts:
- target object observed: yes/no/unknown; evidence step:
- target receptacle observed: yes/no/unknown; evidence step:
- checked locations:
- empty or unhelpful locations:
- inventory state:

Failure diagnosis:
- invalid or malformed actions:
- repeated or wasted actions:
- missing necessary subgoal:

Retry plan:
- search priority:
- first action intention:
- next action intention if target not found:
- action/output format rule:

Rules:
1. If the target object was not observed, write "target not observed yet".
2. Mention only object/receptacle names appearing in observations, actions, or the original task.
3. Give concrete ALFWorld-style action intentions, not generic advice.
4. Keep under 160 words.
\end{lstlisting}

这个 prompt 的意图不是让 summary 更漂亮，而是让它变成 retry policy 的短期工作记忆。对 ALFWorld 来说，最有价值的信息通常不是“我应该更仔细”，而是：
\begin{itemize}
  \item 哪些位置已经查过；
  \item 哪些位置为空或无关；
  \item 目标物是否已经出现；
  \item 下一轮优先查哪里；
  \item 输出格式上要避免什么错误。
\end{itemize}

\section{建议三：增加 Summary Quality Gate}
在 reward 稀疏阶段，不应默认所有 summary 都同等进入训练。建议先记录质量指标，后续再决定是否作为训练权重：
\begin{lstlisting}
faithfulness_pass:
  every mentioned object / receptacle / action is supported by task, obs,
  action history, or final observation.

actionability_pass:
  summary contains concrete search priority or action intention.

format_warning_pass:
  summary tells retry to output exactly one admissible action with
  <think>...</think><action>...</action>.

retry_follow_rate:
  first K traj2 actions follow the search priority in summary.

invalid_delta:
  invalid_count(traj2) - invalid_count(traj1).
\end{lstlisting}

可选训练权重：
\begin{lstlisting}
summary_train_weight =
  1.0 if faithfulness_pass and actionability_pass
  0.3 if actionability_pass but faithfulness uncertain
  0.0 if hallucination or retry invalid rate worsens significantly
\end{lstlisting}

第一版可以只打日志，不过滤；一旦能稳定复现，就把该权重接入 summary row 的 loss weight 或 reward shaping。

\section{建议四：保留 FullTraj 对照，做最小 Ablation}
为了确认“简化 trajectory”本身是否有效，下一批实验建议至少比较五组：
\begin{center}
\small
\begin{tabularx}{\linewidth}{lX}
\toprule
实验组 & 说明 \\
\midrule
SSCA-FullTraj & 当前版本，summary 看完整 prompt / history dump。 \\
SSCA-CompactObsTrace & summary 只看 compact observation / action / feedback trace。 \\
SSCA-CompactObsTrace-NoRawResponse & 不给 raw model response，只给 projected action 和 valid / malformed reason。 \\
SSCA-RandomSummary & 用 random 或 shuffled summary 控制“只是增加上下文长度”的收益。 \\
SSCA-NoSummary & reset 后普通 retry，不提供 summary。 \\
\bottomrule
\end{tabularx}
\end{center}

主要指标：
\begin{lstlisting}
summary_faithfulness_fail_rate
summary_actionability_pass_rate
traj2_invalid_rate
retry_follow_rate
reward2 - reward1
success2 - success1
\end{lstlisting}

在短 smoke 阶段，即使 \texttt{reward2 - reward1} 还没有显著改善，只要 \texttt{CompactObsTrace} 明显降低 hallucination 和 invalid rate，就说明 prompt 方向更健康。

\section{建议五：对应代码改动位置}
建议改动集中在 \path{recipe/SSCA/rollout_loop.py}，不需要碰主 trainer：
\begin{enumerate}
  \item 在 \texttt{\_run\_env\_episode} 中记录 compact step record，而不是只追加字符串型 full prompt history；
  \item 新增 \texttt{\_build\_compact\_summary\_obs(...)}，替代或并列当前 \texttt{\_build\_summary\_obs(...)}；
  \item 新增 malformed reason 提取：missing \texttt{<action>}、missing \texttt{<think>}、contains Chinese、多 action tags、not one admissible action；
  \item 在 summary row 上记录 \texttt{summary\_context\_mode}、\texttt{faithfulness\_pass} 和 \texttt{actionability\_pass}；
  \item 继续记录 \texttt{summary\_train\_weight}，并在 trace viewer 中展示 compact trace、summary gate 和 retry follow rate。
\end{enumerate}

ALFWorld reset 的 \texttt{retry\_same\_seed} 逻辑可以保持不变；这次主要改的是 summary 输入压缩和质量评估。

\section{结论}
当前 SSCA retry smoke 已经验证了链路可行：summary 确实作为 actor response 插入，并和 traj2 共享 retry reward。但逐条审计表明，\texttt{full trajectory} 不是好的 summary 输入形态。它把太多 prompt 级噪声暴露给模型，导致 summary 容易幻觉、泛化、重复计划，甚至让 retry invalid rate 变差。

下一步应优先做两件事：
\begin{enumerate}
  \item 用 \texttt{compact\_obs\_trace} 替代 full prompt dump；
  \item 用 evidence-grounded schema 和 quality gate 约束 summary。
\end{enumerate}

只有当 summary 先满足忠实性和可执行性，再用 retry reward 训练它，SSCA-Suffix 才有机会成为有效的长程 Agent credit source。

\end{document}
